\section{Introduction}
Accurate distance measurement is central to many electrical and electronic engineering applications, ranging from robotics and industrial automation to fluid‑level control and non‑contact gauging. Ultrasonic rangefinders provide a low‑cost, robust method to make such measurements because they are insensitive to ambient light and electromagnetic interference \cite{mi13040520}.  These sensors operate by emitting a short burst of high‑frequency sound and measuring the time elapsed until an echo returns from a target. Non‑contact distance measurement avoids mechanical wear and is safer for applications involving moving parts or hazardous environments.

The ESP32 microcontroller family integrates a dual‑core 32‑bit processor, Wi‑Fi, Bluetooth and a rich set of peripherals including UART, SPI, I²C and ADC interfaces. These modules operate from 3.0–3.6 V, have up to 26 GPIO pins and can clock at up to 240 MHz. Their high computing power and low power consumption make them suitable as the control unit for a digital range finder. An I²C LCD provides a simple user interface to display the measured distance \cite{esp32_datasheet}.

This project aims to design and build a digital range finder that uses ultrasonic time‑of‑flight measurement with the ESP32 and to calibrate the device over a 10–100 cm range. The report describes the physics of ultrasonic ranging, factors affecting measurement accuracy and advanced signal‑processing techniques, then outlines an action plan for implementing and calibrating the instrument.